% GENERATED FILE. Edit canonical lessons, then run npm run generate:books.
% canonical-source-hash: fnv1a64:e6539d5e1e65a1cd
% canonical-lessons: GU-C185-bhavtal-karvo, GU-C185-sodo, GU-C185-ochu, GU-C185-jodi, GU-C185-parakhvu

\chapter{Haggling over a deal}
\label{ch:gu-haggling-over-a-deal}
\hlchaptermodality{185}

\begin{chapteropening}
\textbf{By the end of this chapter:} \emph{I can haggle, ask for a lower rate and judge what I buy.}

Complete the last lesson of chapter 185: I can haggle, ask for a lower rate and judge what I buy.
\end{chapteropening}

\section[\texorpdfstring{\gu{ભાવતાલ} \gu{કરવો} (bhāvtāl karvo)}{ભાવતાલ કરવો (bhāvtāl karvo)}]{\gu{ભાવતાલ} \gu{કરવો} (bhāvtāl karvo) — to haggle, to bargain}

\label{lesson:GU-C185-bhavtal-karvo}



\begin{quote}
Before the new one: say the Gujarati for to measure, then the Gujarati for to ask for.
\end{quote}

\subsection*{You'll want to know: \gu{ભાવતાલ} \gu{કરવો}}
\textbf{\gu{ભાવતાલ} \gu{કરવો}} — \emph{bhāvtāl karvo} — \textquotedblleft{}to haggle, to bargain\textquotedblright{}.

\begin{quote}
\textbf{\gu{તે} \gu{બજારમાં} \gu{ભાવતાલ} \gu{કરે} \gu{છે}.} — \emph{te bajārmā̃ bhāvtāl kare che} — \textquotedblleft{}She haggles in the market.\textquotedblright{}
\end{quote}

Past \textbf{\gu{તેણે} \gu{ભાવતાલ} \gu{કર્યો}}, \emph{teṇe bhāvtāl karyo}.

\begin{cousinweb}
\textbf{\gu{ભાવતાલ} \gu{કરવો}} (\emph{bhāvtāl karvo}) is built from \emph{bhāv}, \textquotedblleft{}rate\textquotedblright{} (CDIAL 9475), paired with \emph{tāl}, and the light verb \emph{karvũ}, \textquotedblleft{}to do\textquotedblright{}. The origin of \emph{tāl} is not settled: CDIAL has no entry for the pair.
\end{cousinweb}

\begin{grammarlens}[title={What you've built}]
One more word for everyday life at A2.
\end{grammarlens}

\subsection*{Your turn}
\begin{itemize}
\raggedright
  \item \emph{Say it:} \emph{bhāvtāl karvo}
  \item \emph{Say it:} \emph{bhāvtāl karvo}, once more
  \item \emph{Say it:} \emph{te bajārmā̃ bhāvtāl kare che} — She haggles in the market.
  \item \emph{Recall:} say the Gujarati for to measure, then the Gujarati for to ask for, then say \emph{bhāvtāl karvo} again
\end{itemize}

\begin{tcolorbox}[breakable,colback=teal!4,colframe=teal!35!black,title={Before you move on}]
What does \gu{ભાવતાલ} \gu{કરવો} mean? (\textquotedblleft{}To haggle, to bargain\textquotedblright{}.) Say it once more.
\end{tcolorbox}
\section[\texorpdfstring{\gu{સોદો} (sodo)}{સોદો (sodo)}]{\gu{સોદો} (sodo) — deal, bargain}

\label{lesson:GU-C185-sodo}



\begin{quote}
Before the new one: say the Gujarati for to ask for, then the Gujarati for to haggle.
\end{quote}

\subsection*{You'll want to know: \gu{સોદો}}
\textbf{\gu{સોદો}} — \emph{sodo} — \textquotedblleft{}deal, bargain\textquotedblright{}.

\begin{quote}
\textbf{\gu{આ} \gu{સોદો} \gu{સારો} \gu{છે}.} — \emph{ā sodo sāro che} — \textquotedblleft{}This is a good deal.\textquotedblright{}
\end{quote}

\begin{cousinweb}
\textbf{\gu{સોદો}} (\emph{sodo}) is a loan, usually traced to Persian \emph{saudā}, \textquotedblleft{}trade, bargain\textquotedblright{}. CDIAL has no entry for it.
\end{cousinweb}

\begin{grammarlens}[title={What you've built}]
One more word for everyday life at A2.
\end{grammarlens}

\subsection*{Your turn}
\begin{itemize}
\raggedright
  \item \emph{Say it:} \emph{sodo}
  \item \emph{Say it:} \emph{sodo}, once more
  \item \emph{Say it:} \emph{ā sodo sāro che} — This is a good deal.
  \item \emph{Recall:} say the Gujarati for to ask for, then the Gujarati for to haggle, then say \emph{sodo} again
\end{itemize}

\begin{tcolorbox}[breakable,colback=teal!4,colframe=teal!35!black,title={Before you move on}]
What does \gu{સોદો} mean? (\textquotedblleft{}Deal, bargain\textquotedblright{}.) Say it once more.
\end{tcolorbox}
\section[\texorpdfstring{\gu{ઓછું} (ochũ)}{ઓછું (ochũ)}]{\gu{ઓછું} (ochũ) — less, too little}

\label{lesson:GU-C185-ochu}



\begin{quote}
Before the new one: say the Gujarati for to haggle, then the Gujarati for deal.
\end{quote}

\subsection*{You'll want to know: \gu{ઓછું}}
\textbf{\gu{ઓછું}} — \emph{ochũ} — \textquotedblleft{}less, too little\textquotedblright{}.

\begin{quote}
\textbf{\gu{આ} \gu{દુકાનમાં} \gu{ભાવ} \gu{ઓછો} \gu{છે}.} — \emph{ā dukānmā̃ bhāv ocho che} — \textquotedblleft{}The rate is lower in this shop.\textquotedblright{}
\end{quote}

\begin{cousinweb}
\textbf{\gu{ઓછું}} (\emph{ochũ}) is a tadbhava. It comes from \emph{*ōccha}, \textquotedblleft{}small, thin, mean\textquotedblright{}, a reconstructed form, and CDIAL glosses the Gujarati word \textquotedblleft{}deficient, inferior\textquotedblright{} (CDIAL 2540).
\end{cousinweb}

\begin{grammarlens}[title={What you've built}]
One more word for everyday life at A2.
\end{grammarlens}

\subsection*{Your turn}
\begin{itemize}
\raggedright
  \item \emph{Say it:} \emph{ochũ}
  \item \emph{Say it:} \emph{ochũ}, once more
  \item \emph{Say it:} \emph{ā dukānmā̃ bhāv ocho che} — The rate is lower in this shop.
  \item \emph{Recall:} say the Gujarati for to haggle, then the Gujarati for deal, then say \emph{ochũ} again
\end{itemize}

\begin{tcolorbox}[breakable,colback=teal!4,colframe=teal!35!black,title={Before you move on}]
What does \gu{ઓછું} mean? (\textquotedblleft{}Less, too little\textquotedblright{}.) Say it once more.
\end{tcolorbox}
\section[\texorpdfstring{\gu{જોડી} (joḍī)}{જોડી (joḍī)}]{\gu{જોડી} (joḍī) — pair}

\label{lesson:GU-C185-jodi}



\begin{quote}
Before the new one: say the Gujarati for deal, then the Gujarati for less.
\end{quote}

\subsection*{You'll want to know: \gu{જોડી}}
\textbf{\gu{જોડી}} — \emph{joḍī} — \textquotedblleft{}pair\textquotedblright{}.

\begin{quote}
\textbf{\gu{આ} \gu{બૂટની} \gu{જોડી} \gu{મોંઘી} \gu{છે}.} — \emph{ā būṭnī joḍī mõghī che} — \textquotedblleft{}This pair of shoes is expensive.\textquotedblright{}
\end{quote}

\begin{cousinweb}
\textbf{\gu{જોડી}} (\emph{joḍī}) is a tadbhava. CDIAL lists it, glossed \textquotedblleft{}pair\textquotedblright{}, under \emph{*yuṭati}, \textquotedblleft{}is joined\textquotedblright{}, a reconstructed form (CDIAL 10496).
\end{cousinweb}

\begin{grammarlens}[title={What you've built}]
One more word for everyday life at A2.
\end{grammarlens}

\subsection*{Your turn}
\begin{itemize}
\raggedright
  \item \emph{Say it:} \emph{joḍī}
  \item \emph{Say it:} \emph{joḍī}, once more
  \item \emph{Say it:} \emph{ā būṭnī joḍī mõghī che} — This pair of shoes is expensive.
  \item \emph{Recall:} say the Gujarati for deal, then the Gujarati for less, then say \emph{joḍī} again
\end{itemize}

\begin{tcolorbox}[breakable,colback=teal!4,colframe=teal!35!black,title={Before you move on}]
What does \gu{જોડી} mean? (\textquotedblleft{}Pair\textquotedblright{}.) Say it once more.
\end{tcolorbox}
\section[\texorpdfstring{\gu{પારખવું} (pārakhvũ)}{પારખવું (pārakhvũ)}]{\gu{પારખવું} (pārakhvũ) — to test, to judge quality}

\label{lesson:GU-C185-parakhvu}



\begin{quote}
Before the new one: say the Gujarati for less, then the Gujarati for pair.
\end{quote}

\subsection*{You'll want to know: \gu{પારખવું}}
\textbf{\gu{પારખવું}} — \emph{pārakhvũ} — \textquotedblleft{}to test, to judge quality\textquotedblright{}.

\begin{quote}
\textbf{\gu{ગ્રાહક} \gu{કેરી} \gu{પારખે} \gu{છે}.} — \emph{grāhak kerī pārakhe che} — \textquotedblleft{}The customer tests the mangoes.\textquotedblright{}
\end{quote}

\begin{cousinweb}
\textbf{\gu{પારખવું}} (\emph{pārakhvũ}) comes from Sanskrit \emph{párīkṣatē}, \textquotedblleft{}observes, examines\textquotedblright{}. CDIAL lists the Gujarati verb as \emph{parakhvũ} and \emph{pārakhvũ}, and marks the long-vowel form as taken from Marathi (CDIAL 7904).
\end{cousinweb}

\begin{grammarlens}[title={What you've built}]
One more word for everyday life at A2.
\end{grammarlens}

\subsection*{Your turn}
\begin{itemize}
\raggedright
  \item \emph{Say it:} \emph{pārakhvũ}
  \item \emph{Say it:} \emph{pārakhvũ}, once more
  \item \emph{Say it:} \emph{grāhak kerī pārakhe che} — The customer tests the mangoes.
  \item \emph{Recall:} say the Gujarati for less, then the Gujarati for pair, then say \emph{pārakhvũ} again
\end{itemize}

\begin{tcolorbox}[breakable,colback=teal!4,colframe=teal!35!black,title={Before you move on}]
What does \gu{પારખવું} mean? (\textquotedblleft{}To test, to judge quality\textquotedblright{}.) Say it once more.
\end{tcolorbox}
